% =====================================================================================
%  ab-ovo-overview.tex
%
%  THIS IS NOT A RESEARCH PAPER, AND IT IS NOT WRITTEN UNDER THE RESEARCH-DOCUMENTATION
%  STANDARD.
%
%  Saying so here, plainly and at the top, is the point of this comment. The document is
%  typeset like a paper because a typeset overview is easier to hand to somebody than a
%  directory of markdown, and a reader who mistakes the form for the genre will expect
%  things this document does not have and must not claim: a research question, a method,
%  data, a result, a threat-to-validity section, a related-work survey, or a finding that
%  can be cited. It has none of those. It is a PROJECT OVERVIEW.
%
%  It also introduces NO FACT its markdown source does not carry. Every statement below is
%  in README.md, docs/architecture/00-ARCHITECTURE.md, docs/adr/ or docs/ux/UI-UX.md, and
%  those files are the source of truth. If this paper and one of them disagree, the
%  markdown is right and this file is stale.
%
%  WHY THAT RULE. A second copy of something that has a source is the defect this estate
%  keeps paying for: the day the two disagree, nothing says which one is the document. The
%  rule reduces the failure to a known direction, which is the most a second rendering can
%  buy. Anything genuinely new belongs in the markdown first, and reaches this file only
%  afterwards.
%
%  BUILD. .github/workflows/build-overview-pdf.yml renders this file on workflow_dispatch
%  and nothing else, and uploads the PDF as a run artifact. THE PDF IS NEVER COMMITTED: a
%  committed PDF is a binary in every diff and a second copy of a document whose source is
%  already here, so the day they disagree nothing says which one is the paper. The .tex is
%  the paper; the workflow renders it.
%
%  The workflow defaults to pdflatex and offers xelatex and lualatex as inputs. This file
%  is deliberately written to compile under pdflatex with a stock TeX Live: standard
%  packages only, and no character outside ASCII. The book it describes sets Polish
%  alongside English; if a future revision of this paper quotes any of it, that revision
%  needs one of the Unicode engines and should say so here.
% =====================================================================================

\documentclass[11pt,a4paper]{article}

\usepackage[T1]{fontenc}
\usepackage[utf8]{inputenc}
\usepackage[a4paper,margin=28mm]{geometry}
\usepackage{booktabs}
\usepackage{enumitem}
\usepackage{parskip}
\usepackage[hidelinks]{hyperref}

\newcommand{\code}[1]{\texttt{#1}}

\title{ab-ovo: a learning platform that encapsulates\\
\emph{Mathematics from Zero for the AI Engineer}}
\author{konradcinkusz}
\date{14 September 2026}

\begin{document}
\maketitle

\begin{abstract}
\noindent
ab-ovo is a learning platform built around a book of 47 programmed-learning programs, in
English and Polish, together with that book's computer exercises. This document is a
project overview: what the system is, what it deliberately is not, what exists in the
repository today, and what is planned. It is not a research paper and reports no result.
Every statement in it is drawn from the repository's markdown documentation, which is the
source of truth; where the two disagree, the markdown is right.
\end{abstract}

\section{What ab-ovo is}

The book it encapsulates is built on a mechanism that paper cannot enforce. A Stroud frame
asks the reader for something \emph{before} it tells them anything, and the next frame opens
with the answer they were supposed to have written. A reader who skims gets nothing from
that structure, and a printed page has no way to notice.

ab-ovo is the software that does notice. A reader reads a frame, commits an answer, and
reveals the next frame, which opens with the answer; where a program carries computer
exercises, they are worked in a lab pane in which Python runs in the reader's own browser.

The product's first requirement constrains everything else in this document: \textbf{the
reader loop works with no account and no backend.} Frames are served with the site as a
versioned content bundle and the exercises run client-side. An account adds synchronisation
of progress between machines, and nothing else.

\section{What it deliberately is not}

The system collects outcomes, and outcomes about learning drift towards becoming outcomes
about learners, because that is the easier number to produce and the one that looks like
progress. The project's answer is an anti-goal stated at the top of the README, above the
features, and enforced by an architectural absence rather than by policy:

\begin{quote}
The instrument measures \textbf{the book}, never the reader.
\end{quote}

The useful question it answers is \emph{where is this book wasting the reader's time}, not
\emph{which reader is worst}. Concretely, and in a form that a search of the code can
falsify: there is no per-reader view and adding one is not planned; there is no endpoint
whose route or query parameter names a person; and there is no table from which a per-reader
score could be built, because the aggregate store holds no reader identity at all. Outcomes
are recorded against a frame in a bundle version, an attempt and a check run.

Two further refusals are recorded as decisions rather than left as preferences. No language
model is called anywhere in the reader loop: the comparison against the next frame is the
teaching, and a model's judgement would carry its own error rate into evidence used to
revise a book. And there is no gamification, because a streak is a reason to move a number
that is not learning.

\section{Status: nothing is deployed}

The repository stands at \emph{builds, tests green, images build}. There is no running
instance, at any address, for anybody. Four \code{fly.toml} files describe a topology that
has never been applied, and neither container image has ever been published.

There is also, deliberately, \textbf{no domain model}: the persistence context declares no
entity and there are no migrations. The mechanism around it --- provider selection, retry,
migration after the listener starts, the health check --- is wired and proven by tests. The
first entities arrive with the first ticket that needs them, because entities invented ahead
of that ticket are code the ticket deletes.

\section{The shape of the repository}

\begin{table}[ht]
\centering
\small
\begin{tabular}{@{}ll@{}}
\toprule
\textbf{Path} & \textbf{What it is} \\
\midrule
\code{src/AbOvo.AppHost}        & The composition root. Development only; not the production topology. \\
\code{src/AbOvo.ServiceDefaults} & The shared kernel: cross-cutting plumbing, under a size ceiling. \\
\code{src/AbOvo.Contracts}      & DTOs that cross a service boundary. \\
\code{src/AbOvo.Api}            & The HTTP service. Owns one logical database. Validates tokens; mints none. \\
\code{tests/AbOvo.Api.Tests}    & Unit and in-memory integration, plus the architecture rules. No container. \\
\code{tests/e2e}                & The Playwright acceptance suite; its own package and lockfile. \\
\code{web/}                     & The pnpm workspace. One Next.js app, and its backend-for-frontend. \\
\code{flyio/}                   & The deployed topology, described and never applied. \\
\code{docs/}                    & Architecture, decisions, UX, diagrams, and this paper. \\
\bottomrule
\end{tabular}
\end{table}

The HTTP surface that exists today is three endpoints: a readiness endpoint reporting the
state of every optional integration, a liveness endpoint carrying only liveness-tagged
checks, and one public endpoint describing the deployment. The authenticated and
administrative route groups are declared in the composition root and carry nothing yet ---
written down before the first endpoint arrives, because a group that does not exist cannot
be seen to be missing.

\section{Architecture, and where it departs from the constitution}

The repository is built against a reference architecture held in a separate standards
repository, and \code{docs/architecture/00-ARCHITECTURE.md} walks its fifteen principles
against this tree --- referencing the constitution rather than copying it, because a copied
principle is a second copy that goes stale.

Three properties are worth naming in an overview. Optional integrations \textbf{degrade
visibly}: a deployment with no credentials at all starts, serves, and names what it does not
have, in three renderings produced from one source. The shared kernel is held to plumbing by
two mechanisms rather than by prose --- a size gate in CI and an architecture test asserting
that it references no service type and declares no persistence model. And no address is
compiled into any artifact: the frontend reads its configuration at request time, so one
image serves every environment.

That document also carries a \textbf{deviation register}, in which every row has a date, a
reason and an exit condition. Five deviations are recorded at the time of writing: dependency
version automation declared and switched off; a size gate counting code lines rather than raw
lines, so that a principle requiring reasoning in comments does not spend the budget of a
principle capping size; Python in the browser, which the standards do not evidence; a
repository name that is a typo and a system name that is not; and one Aspire API marked
evaluation-only, used in a development-only file.

A deviation with no exit condition is a drift rather than a decision, which is why the
column exists.

\section{The plan}

Five phases, ranked in \code{docs/ux/UI-UX.md} so that the first delivery session picks up a
decision rather than re-deriving one.

\begin{enumerate}[leftmargin=2em,itemsep=2pt]
  \item \textbf{The lab pane.} The exercises, in the browser. The largest piece that needs
        neither a backend nor the content bundle, which is why it is first.
  \item \textbf{The content schema and the frame view.} Its first half proceeds against a
        fixture; its second half is \emph{blocked} on the book publishing a content bundle,
        which is the plan's only external dependency.
  \item \textbf{Progress and accounts.} Local progress first; an account synchronises it.
  \item \textbf{The instrument.} Opt-in, versioned consent; outcomes against artifacts; every
        rate carrying its interval.
  \item \textbf{The open-source edition.} A one-time, irreversible gate, ordered by
        irreversibility, in which a full history scan for secrets blocks every other item.
\end{enumerate}

\section{Where the documents are}

\begin{description}[leftmargin=0pt,itemsep=1pt]
  \item[\code{README.md}] What the system is, how to run it in one command, and the
        anti-goals, stated before the features.
  \item[\code{docs/architecture/00-ARCHITECTURE.md}] The principle-by-principle walk, the
        deviation register, and the known gaps.
  \item[\code{docs/adr/}] The decisions, one file each, with what they cost.
  \item[\code{docs/ux/UI-UX.md}] The screens as scaffolded, and the ranked backlog.
  \item[\code{docs/diagrams/}] The system topology and the reader loop.
  \item[\code{AGENTS.md}, \code{CONTRIBUTING.md}, \code{SECURITY.md}] Working in the
        repository, and what to do when a secret lands in history.
\end{description}

\vspace{1em}
\noindent\rule{\linewidth}{0.4pt}

\small
\noindent
This document records no measurement of its own and makes no claim that is not already made
in the repository's markdown. It is rendered on request by a manual workflow, and the
resulting PDF is not committed.

\end{document}
